\section{Conclusion}
Les résultats mettent en évidence des différences importantes entre les plots TS et les plots T. Sur le plan du sol, les plots TS ont conservé une structure meuble et une meilleure infiltration, tandis que les plots T ont présenté un durcissement, une fissuration, un ruissellement préférentiel et des efflorescences blanchâtres, ce qui est cohérent avec la salinité et la sodicité observées au départ. 
Les plots TS ont aussi montré une meilleure survie et une croissance plus importante (gombo TS : 41,3 cm contre 21 cm pour le témoin), qui s’est traduite par une production effective : récolte régulière pour le gombo, 4,8~kg sur $2\,\mathrm{m^2}$ pour le navet et récolte de laitue malgré les pertes, alors que les plots T ont produit très peu, voire pas du tout. Le pilotage par capteur a en outre réduit la consommation d’eau de 66~\% à 86~\% selon les cultures.

Ces observations montrent une association entre l’amendement du sol, l’amélioration de son état physique et les performances des plots TS. Certains phénomènes, notamment la coloration noirâtre de l’eau après arrosage, nécessitent toutefois des analyses complémentaires. Ces résultats constituent une base pour les recommandations et perspectives présentées dans la suite du mémoire.